% ===========================================================================
\chapter{Modelo de Domínio e Entidades de Negócio (Domain Layer)}
\label{cap_modelo_dominio}
% ===========================================================================

A camada de domínio concentra a inteligência de negócio, as regras clínicas da metodologia IHI-GTT e as restrições invariantes do SIGEA-GTT. Projetada sob os preceitos do \textit{Domain-Driven Design} (DDD) \cite{evans2003}, esta camada é livre de dependências de infraestrutura externa ou protocolos de transporte, sendo composta por Entidades JPA fortemente tipadas, Enumerações de controle de acesso e Objetos de Valor (\textit{Value Objects}) que modelam documentos clínicos semiestruturados.

% ===========================================================================
\section{Diagrama de Classes de Domínio}
\label{sec_diagrama_classes_dominio}
% ===========================================================================

A Figura \ref{fig_diagrama_classes_dominio} apresenta o Diagrama Geral de Classes de Domínio do SIGEA-GTT, renderizado em notação gráfica nativa TikZ com fidelidade à especificação OMG UML 2.5.1:

\begin{figure}[!ht]
\centering
\caption{Diagrama de Classes de Domínio do SIGEA-GTT (Entidades JPA)}
\label{fig_diagrama_classes_dominio}
\resizebox{0.96\textwidth}{!}{%
\begin{tikzpicture}[
    class_box/.style={
        rectangle split,
        rectangle split parts=3,
        rectangle split part fill={clinical-100!30, white, white},
        draw=clinical-700,
        thick,
        rounded corners=2pt,
        font=\tiny,
        minimum width=4.8cm,
        text width=4.6cm,
        align=left
    },
    enum_box/.style={
        rectangle split,
        rectangle split parts=2,
        rectangle split part fill={alert-amber!20, white},
        draw=alert-amber!90!black,
        thick,
        rounded corners=2pt,
        font=\tiny,
        minimum width=3.0cm,
        text width=2.8cm,
        align=left
    },
    assoc/.style={
        draw=clinical-900,
        thick,
        -Latex
    },
    comp/.style={
        draw=clinical-900,
        thick,
        {Diamond[fill=clinical-900]}-Latex
    }
]

    % ====================================================
    % ENUM USER_ROLE
    % ====================================================
    \node[enum_box] (user_role) at (-4.5, 5.5) {
        \centering\textbf{\scriptsize <<enumeration>>\\UserRole}
        \nodepart{two}
        + ADMIN\\
        + PROFESSOR\\
        + STUDENT
    };

    % ====================================================
    % CLASSE USER
    % ====================================================
    \node[class_box] (user) at (0.5, 5.5) {
        \centering\textbf{\scriptsize <<Entity>>\\User}
        \nodepart{two}
        - id : UUID\\
        - email : String\\
        - fullName : String\\
        - passwordHash : String\\
        - role : UserRole\\
        - registrationNumber : String\\
        - isActive : Boolean\\
        - mustChangePassword : Boolean\\
        - createdAt : Instant\\
        - updatedAt : Instant
        \nodepart{three}
        + isAdmin() : boolean\\
        + isProfessor() : boolean\\
        + isStudent() : boolean
    };

    % ====================================================
    % CLASSE ACADEMIC_CLASS
    % ====================================================
    \node[class_box] (class) at (7.5, 5.5) {
        \centering\textbf{\scriptsize <<Entity>>\\AcademicClass}
        \nodepart{two}
        - id : UUID\\
        - subjectName : String\\
        - classCode : String\\
        - academicPeriod : String\\
        - professor : User\\
        - students : Set<User>\\
        - isClosed : Boolean\\
        - createdAt : Instant
        \nodepart{three}
        + getFormattedName() : String\\
        + enrollStudent(User) : void\\
        + removeStudent(User) : void\\
        + closeClass() : void
    };

    % ====================================================
    % CLASSE ACTIVITY
    % ====================================================
    \node[class_box] (activity) at (7.5, -0.5) {
        \centering\textbf{\scriptsize <<Entity>>\\Activity}
        \nodepart{two}
        - id : UUID\\
        - academicClass : AcademicClass\\
        - title : String\\
        - description : String\\
        - clinicalCaseData : ClinicalCaseData\\
        - deadline : Instant\\
        - createdAt : Instant
        \nodepart{three}
        + isExpired() : boolean\\
        + isEligibleForSubmission() : boolean
    };

    % ====================================================
    % CLASSE ACTIVITY_SUBMISSION
    % ====================================================
    \node[class_box] (submission) at (0.5, -0.5) {
        \centering\textbf{\scriptsize <<Entity>>\\ActivitySubmission}
        \nodepart{two}
        - id : UUID\\
        - activity : Activity\\
        - student : User\\
        - identifiedTriggers : List<IdentifiedTriggerData>\\
        - qualityToolsData : QualityToolsData\\
        - submissionDate : Instant\\
        - grade : BigDecimal\\
        - professorFeedback : String\\
        - gradedAt : Instant
        \nodepart{three}
        + isGraded() : boolean\\
        + assignGrade(BigDecimal, String) : void
    };

    % ====================================================
    % CLASSE GTT_MODULE
    % ====================================================
    \node[class_box] (module) at (-4.5, -6.5) {
        \centering\textbf{\scriptsize <<Entity>>\\GttModule}
        \nodepart{two}
        - id : UUID\\
        - code : String\\
        - name : String\\
        - description : String\\
        - isActive : Boolean\\
        - createdAt : Instant
        \nodepart{three}
        + getDisplayName() : String\\
        + activate() : void\\
        + deactivate() : void
    };

    % ====================================================
    % CLASSE GTT_TRIGGER
    % ====================================================
    \node[class_box] (trigger) at (1.5, -6.5) {
        \centering\textbf{\scriptsize <<Entity>>\\GttTrigger}
        \nodepart{two}
        - id : UUID\\
        - module : GttModule\\
        - code : String\\
        - name : String\\
        - description : String\\
        - isActive : Boolean\\
        - createdAt : Instant
        \nodepart{three}
        + getFullCode() : String
    };

    % ====================================================
    % CLASSE HARM_SEVERITY
    % ====================================================
    \node[class_box] (severity) at (7.5, -6.5) {
        \centering\textbf{\scriptsize <<Entity>>\\HarmSeverity}
        \nodepart{two}
        - id : UUID\\
        - categoryLetter : String\\
        - name : String\\
        - description : String\\
        - isHarm : Boolean\\
        - isActive : Boolean
        \nodepart{three}
        + isHarmCategory() : boolean
    };

    % ====================================================
    % RELACIONAMENTOS E NAVEGABILIDADE
    % ====================================================
    % User -> UserRole
    \draw[assoc] (user) -- node[above, font=\tiny] {role} (user_role);

    % AcademicClass -> User (Professor)
    \draw[assoc] (class) -- node[above, font=\tiny] {1 \quad professor \quad 1} (user);

    % AcademicClass -> User (Students)
    \draw[assoc] (class.north) to[bend right=25] node[above, font=\tiny] {* \quad students \quad *} (user.north);

    % AcademicClass -> Activity (1:N)
    \draw[comp] (class) -- node[right, font=\tiny] {1 \quad 1..*} (activity);

    % Activity -> ActivitySubmission (1:N)
    \draw[comp] (activity) -- node[above, font=\tiny] {1 \quad *} (submission);

    % User -> ActivitySubmission (Student)
    \draw[assoc] (user) -- node[right, font=\tiny] {1 \quad student \quad *} (submission);

    % GttModule -> GttTrigger (1:N)
    \draw[comp] (module) -- node[above, font=\tiny] {1 \quad 1..*} (trigger);

\end{tikzpicture}
}
\fonte{Elaborado pelo autor (2026).}
\end{figure}

% ===========================================================================
\section{Detalhamento das Entidades JPA de Negócio}
\label{sec_detalhamento_entidades}
% ===========================================================================

A seguir são descritas as especificações formais de atributos, tipos, restrições e comportamentos das entidades mapeadas:

% ---------------------------------------------------------------------------
\subsection{Entidade: \texttt{User}}
A classe \texttt{User} modela os atores humanos cadastrados no sistema, implementando a segregação de perfis de acordo com o padrão RBAC:
\begin{itemize}
    \item \texttt{id (UUID)}: Chave primária sintética gerada por algoritmo de alta dispersão (UUID v4), impedindo enumeração sequencial maliciosa;
    \item \texttt{email (String)}: Endereço institucional obrigatório com constraint de unicidade e validação de domínio estrito (\texttt{@academico.ufs.br});
    \item \texttt{fullName (String)}: Nome completo do usuário para exibição em listas e relatórios;
    \item \texttt{passwordHash (String)}: Hash unidirecional salgado gerado pelo algoritmo criptográfico BCrypt com fator de custo 10;
    \item \texttt{role (UserRole)}: Enumeração que categoriza o papel operacional do usuário (\texttt{ADMIN}, \texttt{PROFESSOR} ou \texttt{STUDENT});
    \item \texttt{registrationNumber (String)}: Matrícula institucional discente. Obrigatória estritamente quando \texttt{role == UserRole.STUDENT} e obrigatoriamente nula para docentes e administradores;
    \item \texttt{isActive (Boolean)}: Sinalizador booleano de habilitação de acesso, permitindo inativação lógica sem perda de histórico acadêmico;
    \item \texttt{mustChangePassword (Boolean)}: Sinalizador de obrigatoriedade de redefinição de senha provisória gerada automaticamente.
\end{itemize}

% ---------------------------------------------------------------------------
\subsection{Entidade: \texttt{AcademicClass}}
A classe \texttt{AcademicClass} encapsula o agrupamento semestral de discentes sob a supervisão pedagógica de um docente titular:
\begin{itemize}
    \item \texttt{subjectName (String)}: Denominação formal da disciplina curricular (ex.: \textit{"Enfermagem na Saúde do Adulto I"});
    \item \texttt{classCode (String)}: Identificador de turma semestral (ex.: \textit{"T01"});
    \item \texttt{academicPeriod (String)}: Período letivo no formato canônico da UFS (ex.: \textit{"2025.2"});
    \item \texttt{professor (User)}: Referência obrigatória (\texttt{@NotNull}) ao docente titular da turma;
    \item \texttt{students (Set<User>)}: Coleção de discentes matriculados, mapeada pela tabela associativa \texttt{class\_students};
    \item \texttt{isClosed (Boolean)}: Sinalizador de encerramento acadêmico. Quando verdadeiro, veda novas matrículas e submissões tardias;
    \item Método \texttt{getFormattedName()}: Retorna a denominação concatenada \texttt{"Disciplina - Turma - Período"}.
\end{itemize}

% ---------------------------------------------------------------------------
\subsection{Entidades: \texttt{GttModule}, \texttt{GttTrigger} e \texttt{HarmSeverity}}
Estas entidades formam o núcleo taxonômico da metodologia \textit{Global Trigger Tool}:
\begin{itemize}
    \item \texttt{GttModule}: Representa os seis módulos temáticos recomendados pelo IHI \cite{ihi2009}: Cuidados Gerais (\texttt{C}), Cirúrgico (\texttt{S}), Medicamentoso (\texttt{M}), Terapia Intensiva (\texttt{I}), Perinatal (\texttt{P}) e Emergência (\texttt{E});
    \item \texttt{GttTrigger}: Representa os cinquenta e três gatilhos clínicos padronizados. Possui código alfanumérico único (ex.: \texttt{"M03"}) e chave estrangeira mandatória para seu respectivo módulo (\texttt{ON DELETE RESTRICT});
    \item \texttt{HarmSeverity}: Representa as nove categorias de dano do NCC MERP \cite{nccmerp2001} (letras \texttt{'A'} até \texttt{'I'}). O atributo booleano \texttt{isHarm} diferencia incidentes sem dano (\texttt{false} para categorias A até D) de eventos adversos com dano real (\texttt{true} para categorias E até I).
\end{itemize}

% ---------------------------------------------------------------------------
\subsection{Entidades: \texttt{Activity} e \texttt{ActivitySubmission}}
Estruturam o ciclo de realização e avaliação de auditorias clínicas:
\begin{itemize}
    \item \texttt{Activity}: Modela uma atividade de auditoria designada para uma turma. Contém título, instruções, prazo final (\texttt{deadline}) e o prontuário de alta fidelidade persistido na coluna estruturada \texttt{clinicalCaseData};
    \item \texttt{ActivitySubmission}: Armazena a resolução entregue pelo discente. Agrupa a lista de gatilhos auditados (\texttt{identifiedTriggers}), os artefatos de qualidade (\texttt{qualityToolsData}), a data de envio, a nota atribuída pelo docente (\texttt{grade}) e o parecer pedagógico formativo registrado.
\end{itemize}

% ===========================================================================
\section{Objetos de Valor e Modelagem de Dados Clínicos em JSONB}
\label{sec_value_objects_json}
% ===========================================================================

Para acomodar a dinâmica e a heterogeneidade das informações médicas sem incorrer na rigidez excessiva do modelo relacional puro, o sistema adota Objetos de Valor (\textit{Value Objects}) que são serializados em colunas \texttt{JSONB} do PostgreSQL:

\begin{enumerate}
    \item \textbf{\texttt{ClinicalCaseData}}: Agrega os nós do prontuário simulado:
    \begin{itemize}
        \item \texttt{PatientData}: Dados sociodemográficos fictícios (nome, idade, sexo, leito);
        \item \texttt{AdmissionData}: Data de internação, queixa principal e hipótese diagnóstica;
        \item \texttt{List<EvolutionNoteData>}: Evoluções médicas diárias sequenciais;
        \item \texttt{List<PrescriptionData>}: Prescrições farmacológicas ativas e suspensas com dose, via e frequência;
        \item \texttt{List<LabExamData>}: Painéis laboratoriais com resultados numéricos, unidades e faixas de referência;
        \item \texttt{List<ProcedureData>}: Descrições cirúrgicas e procedimentos invasivos.
    \end{itemize}

    \item \textbf{\texttt{IdentifiedTriggerData}}: Registra os rastros clínicos detectados pelo discente:
    \begin{itemize}
        \item Código e nome do gatilho rastreado;
        \item Seção do prontuário onde o rastro foi localizado;
        \item Trecho textual citado da evidência (\texttt{quotedEvidence});
        \item Justificativa clínica redigida pelo auditor (\texttt{clinicalJustification});
        \item Categoria de severidade atribuída na escala NCC MERP (\texttt{'A'} a \texttt{'I'}).
    \end{itemize}

    \item \textbf{\texttt{QualityToolsData}}: Consolida as quatro ferramentas de melhoria contínua:
    \begin{itemize}
        \item \texttt{IshikawaData}: Efeito central e causas nos 6 vetores causais (Método, Máquina, Material, Mão de Obra, Medida, Meio Ambiente);
        \item \texttt{List<GutItemData>}: Problemas priorizados com notas de 1 a 5 para Gravidade, Urgência e Tendência, e score $GUT$;
        \item \texttt{List<FiveWTwoHItemData>}: Ações corretivas detalhadas nas dimensões do método 5W2H;
        \item \texttt{PdcaData}: Metas e rotinas nos quatro quadrantes do ciclo PDCA.
    \end{itemize}
\end{enumerate}
